\begin{proof}[Proof of \cref{obj:thm:full-experiment-lower}]
Let \(\underline c_{\mathrm{one}}>0\) be the universal constant supplied by
\cref{obj:lem:one-cell-testing-family}. Define the universal constants
\[
\eta=e^{-32},\qquad
N_{\mathrm{act}}
=\max\left\{1,e^2,4\eta,4096\eta,
e^{(32+\log 8)/28}\right\},
\]
\[
\underline c_{\mathrm{act}}
=\min\left\{
\frac{1}{256N_{\mathrm{act}}},
\frac{1}{256\cdot2048^2},
\frac{\eta^2}{16384}
\right\},
\qquad
\underline c_*=\min\{\underline c_{\mathrm{one}},
\underline c_{\mathrm{act}}\},
\qquad
\underline c=\frac{\underline c_*}{4}.
\]
These constants are positive and independent of \(n,d,q\).

Fix parameters satisfying the hypotheses, and define
\[
N=nq,\qquad \ell=\log(e+N),\qquad
u_{\mathrm{act}}=\frac{d}{N\ell},\qquad
a_{\mathrm{eff}}=N^{-1},\qquad
b_{\mathrm{dim}}=u_{\mathrm{act}}^2.
\]
Then
\[
N>0,\qquad
\ell\ge\log e=1,\qquad
a_{\mathrm{eff}}\ge0,\qquad b_{\mathrm{dim}}\ge0.
\]
By \cref{obj:def:frontier-rate},
\[
r_{n,d,q}=\min\{1,a_{\mathrm{eff}}+b_{\mathrm{dim}}\}.
\]

\begin{enumerate}
\item \emph{The one-cell contribution.}
By \cref{obj:lem:one-cell-testing-family},
\[
\underline c_{\mathrm{one}}\min\{1,a_{\mathrm{eff}}\}
\le \mathfrak R_{n,d,q}.
\]
We next establish the complementary bound
\[
\sqrt N\le d
\quad\Longrightarrow\quad
\underline c_{\mathrm{act}}\min\{1,b_{\mathrm{dim}}\}
\le \mathfrak R_{n,d,q}.
\]
% lean: CausalSmith.Stat.MarRareqLogfrontier.one_cell_testing_family

\item \emph{A numerical one-cell estimate for constant calibration.}
We record explicitly the numerical estimate used below:
\[
\frac1{256}\min\{1,N^{-1}\}\le\mathfrak R_{n,d,q}.
\]
For \(\mu_{\mathrm{one}}\in[1/4,3/4]\), define \(P_{\mathrm{one},\mu_{\mathrm{one}}}\) by
\[
X=0,\qquad S(0)=S(1)=Y(0)=0,\qquad
(A,Y(1),R)\sim
\operatorname{Bernoulli}(1/2)\otimes
\operatorname{Bernoulli}(\mu_{\mathrm{one}})\otimes
\operatorname{Bernoulli}(q),
\]
\[
S=0,\qquad Y=AY(1).
\]
Treatment is independent of the potential variables, and arrival is independent
of the outcome conditional on the observed cell. Every occupied cell has arrival
probability \(q\). Consequently,
\[
P_{\mathrm{one},\mu_{\mathrm{one}}}\in\mathcal M_{n,d,q},
\qquad
\tau(P_{\mathrm{one},\mu_{\mathrm{one}}})=\mu_{\mathrm{one}}.
\]
Write
\[
Q_{\mathrm{one},\mu_{\mathrm{one}}}=\mathcal L_{P_{\mathrm{one},\mu_{\mathrm{one}}}}(O_1),
\qquad
\delta_{\mathrm{one}}
=\min\left\{\frac14,\frac{1}{4\sqrt N}\right\}.
\]
Thus \(0<\delta_{\mathrm{one}}\le1/4\), and the two target values
\(1/2\) and \(1/2+\delta_{\mathrm{one}}\) belong to \([1/4,3/4]\).
Considering the two branches of this minimum gives
\[
\delta_{\mathrm{one}}^2
\ge \frac1{16}\min\{1,N^{-1}\},
\qquad
16N\delta_{\mathrm{one}}^2\le1.
\]

For probability laws \(Q'\ll Q\), use the convention
\[
\chi^2(Q'\Vert Q)
=\int\left(\frac{dQ'}{dQ}-1\right)^2\,dQ.
\]
The two one-record laws differ precisely at the treated, arrived records with
outcome zero or one. Each of these atoms has probability \(q/4\) under
\(Q_{\mathrm{one},1/2}\), and its probability changes in absolute value by
\(q\delta_{\mathrm{one}}/2\). Summing their contributions gives
\[
\chi^2\!\left(
Q_{\mathrm{one},1/2+\delta_{\mathrm{one}}}
\Vert Q_{\mathrm{one},1/2}\right)
=2q\delta_{\mathrm{one}}^2
\le\frac{1}{8n}.
\]
Independence of the records factors the second moment of the likelihood ratio,
so
\[
1+\chi^2\!\left(
Q_{\mathrm{one},1/2+\delta_{\mathrm{one}}}^{\otimes n}
\Vert Q_{\mathrm{one},1/2}^{\otimes n}\right)
=
\left[
1+\chi^2\!\left(
Q_{\mathrm{one},1/2+\delta_{\mathrm{one}}}
\Vert Q_{\mathrm{one},1/2}\right)
\right]^n
\le e^{1/8}.
\]
Since \(e^{1/8}\le e^{1/2}<\sqrt3<2\), Cauchy--Schwarz applied to the
likelihood ratio yields
\[
\operatorname{TV}\!\left(
Q_{\mathrm{one},1/2+\delta_{\mathrm{one}}}^{\otimes n},
Q_{\mathrm{one},1/2}^{\otimes n}\right)
\le
\frac12\sqrt{
\chi^2\!\left(
Q_{\mathrm{one},1/2+\delta_{\mathrm{one}}}^{\otimes n}
\Vert Q_{\mathrm{one},1/2}^{\otimes n}\right)}
\le\frac12.
\]

For every \(\mathsf T\in\mathcal T_n\), define its barycenter on the
observed-sample space by
\[
b_{\mathsf T}(o)=\int_{[-1,1]}h\,\mathsf T(o,dh).
\]
By \cref{obj:lem:randomized-kernel-barycenter},
\[
b_{\mathsf T}(o)\in[-1,1],
\qquad
E_P[(b_{\mathsf T}(\mathbf O_n)-\tau(P))^2]
\le E_P[(\mathsf T-\tau(P))^2].
\]
Define the nearer-target test
\[
\psi_{\mathrm{one}}(o)
=
\mathbf1\left\{
\left|b_{\mathsf T}(o)-\left(\frac12+\delta_{\mathrm{one}}\right)\right|
\le
\left|b_{\mathsf T}(o)-\frac12\right|
\right\}.
\]
An error under either target entails absolute estimation error at least
\(\delta_{\mathrm{one}}/2\). The testing inequality in
\cref{obj:lem:randomized-kernel-tv-comparison}, followed by the barycenter
risk comparison, therefore gives
\[
\begin{aligned}
&\max\left\{
E_{P_{\mathrm{one},1/2}}\!\left[(\mathsf T-\tfrac12)^2\right],
E_{P_{\mathrm{one},1/2+\delta_{\mathrm{one}}}}\!
\left[(\mathsf T-\tfrac12-\delta_{\mathrm{one}})^2\right]
\right\}\\
&\qquad\ge
\frac{\delta_{\mathrm{one}}^2}{8}
\left[
1-\operatorname{TV}\!\left(
Q_{\mathrm{one},1/2+\delta_{\mathrm{one}}}^{\otimes n},
Q_{\mathrm{one},1/2}^{\otimes n}\right)
\right]
\ge\frac{\delta_{\mathrm{one}}^2}{16}
\ge\frac1{256}\min\{1,N^{-1}\}.
\end{aligned}
\]
Both laws belong to the model class. Taking the supremum over that class and
then the infimum over \(\mathsf T\) proves the displayed numerical estimate.
% lean: CausalSmith.Stat.MarRareqLogfrontier.oneCell_minimax_risk_floors

\item \emph{Large effective size and many rare cells.}
To prove the complementary bound, suppose \(\sqrt N\le d\).
Use the construction parameters of
\cref{obj:def:activation-lower-handle}:
\[
K=\lceil\ell\rceil,\qquad H=K^2,\qquad
b_0=\frac{\eta}{N\ell},\qquad
J=\min\left\{d-1,\left\lfloor\frac1{2b_0}\right\rfloor\right\}.
\]
Here \(b_0>0\), \(0\le J\le d-1\), and
\[
Jb_0\le\frac12.
\]
The ceiling bound and \(\ell\ge1\) give
\[
\ell\le K<\ell+1\le2\ell,\qquad H\ge1.
\]
The rare-arrival restriction then implies
\[
0<qH\le4q\ell^2\le\frac1{16}.
\]

First suppose \(N_{\mathrm{act}}\le N\). In particular,
\[
N\ge1,\qquad N\ge e^2,\qquad
N\ge4\eta,\qquad N\ge4096\eta,\qquad
N\ge e^{(32+\log8)/28}.
\]
The capacity bound is
\[
2048(2b_0)=\frac{4096\eta}{N\ell}\le1,
\qquad
2048\le\frac1{2b_0}.
\]
Within this branch, first suppose \(J\ge2048\). Since \(N\ge e^2\),
\[
\ell=\log(e+N)\ge2,\qquad K\ge2.
\]
By \cref{obj:lem:moment-matched-reciprocal-priors}, there are finitely
supported probability laws \(\Pi_0,\Pi_1\) on \([1,H]\) satisfying
\[
\int z^v\,d\Pi_0(z)=\int z^v\,d\Pi_1(z)
\qquad(0\le v\le K,\ v\in\mathbb N).
\]
If the reciprocal mean under \(\Pi_0\) exceeds that under \(\Pi_1\), interchange
the priors; otherwise retain their order. Their separation then reads
\[
\int z^{-1}\,d\Pi_1(z)-\int z^{-1}\,d\Pi_0(z)
\ge\frac1{12}.
\]

For each \(\mathbf z=(z_x)_{x=0}^{d-1}\in[1,H]^d\), define a full-data law
\(P_{\mathrm{act},\mathbf z}\) as follows. Its baseline probabilities are
\[
p_{\mathrm{act},x}=
\begin{cases}
b_0,&0\le x<J,\\
1-Jb_0,&x=J,\\
0,&J<x<d.
\end{cases}
\]
Set \(S(0)=S(1)=Y(0)=0\), draw \(A\sim\operatorname{Bernoulli}(1/2)\)
independently of the baseline and potential variables, and impose \(S=0\) and
\(Y=AY(1)\). Conditional on \(X=x\), draw \(Y(1)\) and \(R\) independently with
\[
Y(1)\mid X=x\sim
\begin{cases}
\operatorname{Bernoulli}(z_x^{-1}),&x<J,\\
\operatorname{Bernoulli}(0),&x\ge J,
\end{cases}
\qquad
R\mid X=x\sim
\begin{cases}
\operatorname{Bernoulli}(qz_x),&x<J,\\
\operatorname{Bernoulli}(qH),&x\ge J.
\end{cases}
\]
The baseline masses sum to one and are nonnegative. Since \(1\le z_x\le H\),
the occupied-cell arrival probabilities belong to \([q,1]\). The specified
independences give randomized assignment and conditional independence of
arrival and outcome, and the definitions of \(S,Y\) give consistency.
With independent sampling from this fixed law, all model conditions hold:
\[
P_{\mathrm{act},\mathbf z}\in\mathcal M_{n,d,q},
\qquad
\tau_{\mathrm{act}}(\mathbf z)
:=\tau(P_{\mathrm{act},\mathbf z})
=b_0\sum_{x=0}^{J-1}z_x^{-1}.
\]

For \(\iota\in\{0,1\}\), draw
\[
\mathbf Z=(Z_x)_{x=0}^{d-1}\sim\Pi_\iota^{\otimes d}
\]
once for the entire sample. Define the observed mixture law
\[
Q_{\mathrm{act},\iota}
=\int_{[1,H]^d}
\mathcal L_{P_{\mathrm{act},\mathbf z}}(\mathbf O_n)\,
d\Pi_\iota^{\otimes d}(\mathbf z),
\]
and let \(\widetilde Q_{\mathrm{act},\iota}\) be the corresponding augmented
law constructed in \cref{obj:def:activation-lower-handle}. Thus
\[
\widetilde Q_{\mathrm{act},\iota}
=
\int_{[1,H]^d}
\mathcal L\!\left(
(\mathsf B_{\mathrm{act},t},O_t)_{t=1}^{n}
\,\middle|\,\mathbf Z=\mathbf z\right)
d\Pi_\iota^{\otimes d}(\mathbf z).
\]
The activation flags have probability \(qH\), independently of the latent
vector and the baseline and treatment pattern.

At the record level, forgetting the flag gives the following probabilities
for the three marks \((R,RY)=(1,1),(1,0),(0,0)\):
\[
\begin{array}{c|ccc}
& (1,1)&(1,0)&(0,0)\\ \hline
A=1,\ X=x<J&q&q(z_x-1)&1-qz_x\\
A=0,\ X=x<J&0&qz_x&1-qz_x\\
X\ge J&0&qH&1-qH
\end{array}
\]
These are exactly the probabilities induced by \(P_{\mathrm{act},\mathbf z}\).
Define the forgetting map by
\[
\operatorname{pr}_{\mathrm{obs}}
\bigl((\mathsf B_{\mathrm{act},t},O_t)_{t=1}^{n}\bigr)
=(O_t)_{t=1}^{n}.
\]
Conditional independence of the records and integration over the common
latent vector therefore give
\[
Q_{\mathrm{act},\iota}
=\widetilde Q_{\mathrm{act},\iota}
\circ\operatorname{pr}_{\mathrm{obs}}^{-1}.
\]
% lean: CausalSmith.Stat.MarRareqLogfrontier.interval_ordered_reciprocal_priors

\item \emph{Equality on the moment-matching event.}
Continue in the branch \(N\ge N_{\mathrm{act}}\), \(J\ge2048\).
For an augmented sample, define the activated count in cell \(x\) and the
cutoff event by
\[
V_x=\sum_{t=1}^{n}
\mathbf1\{X_t=x,\ \mathsf B_{\mathrm{act},t}=1\},
\qquad
\mathcal G=\{V_x\le K\text{ for every }0\le x<J\}.
\]

Fix a possible augmented sample \(\mathbf o_{\mathrm{aug}}\), and denote its
coordinates by the corresponding record symbols. Its baseline, treatment,
surrogate, and flag pattern has the common probability
\[
\varpi_{\mathbf o_{\mathrm{aug}}}
=
\mathbf1\{S_t=0\text{ for every }t\}
\prod_{t=1}^{n}
\frac{p_{\mathrm{act},X_t}}2
(qH)^{\mathsf B_{\mathrm{act},t}}
(1-qH)^{1-\mathsf B_{\mathrm{act},t}}.
\]
For each cell \(x\), define \(\mathcal P_{\mathbf o_{\mathrm{aug}},x}(z)\)
as the product, over records with \(X_t=x\), of their conditional mark
probabilities from \cref{obj:def:activation-lower-handle}:
\[
\mathcal P_{\mathbf o_{\mathrm{aug}},x}(z)
=
\prod_{t:X_t=x}
\Pr\!\left(
(R,RY)=(R_t,R_tY_t)
\,\middle|\,
X=x,\ A=A_t,\ \mathsf B_{\mathrm{act}}=\mathsf B_{\mathrm{act},t},
\ Z_x=z
\right).
\]
Empty products equal one. The row entries extend these expressions to
polynomials on \(\mathbb R\); their probability interpretation applies on
\([1,H]\). Impossible marks have weight zero. For a rare cell, an inactive
factor is constant and an activated factor is affine in \(z\). For every
other cell the factors are constant. Hence
\[
\deg\mathcal P_{\mathbf o_{\mathrm{aug}},x}\le V_x
\quad(x<J),\qquad
\deg\mathcal P_{\mathbf o_{\mathrm{aug}},x}\le0
\quad(x\ge J),
\]
where a zero polynomial satisfies each stated degree bound.

On \(\mathcal G\), define the coefficients
\(\beta_{\mathbf o_{\mathrm{aug}},x,v}\) by the polynomial expansion
\[
\mathcal P_{\mathbf o_{\mathrm{aug}},x}(z)
=\sum_{v=0}^{K}
\beta_{\mathbf o_{\mathrm{aug}},x,v}z^v
\qquad(z\in\mathbb R),
\]
padding with zero coefficients when necessary. Product-prior independence
then gives
\[
\widetilde Q_{\mathrm{act},\iota}\{\mathbf o_{\mathrm{aug}}\}
=
\varpi_{\mathbf o_{\mathrm{aug}}}
\prod_{x=0}^{d-1}
\left[
\sum_{v=0}^{K}
\beta_{\mathbf o_{\mathrm{aug}},x,v}
\int z^v\,d\Pi_\iota(z)
\right]
\qquad(\mathbf o_{\mathrm{aug}}\in\mathcal G).
\]
The matched moments make the right-hand side identical for
\(\iota=0,1\). Thus
\[
\widetilde Q_{\mathrm{act},0}\{\mathbf o_{\mathrm{aug}}\}
=
\widetilde Q_{\mathrm{act},1}\{\mathbf o_{\mathrm{aug}}\}
\qquad(\mathbf o_{\mathrm{aug}}\in\mathcal G).
\]
% lean: CausalSmith.Stat.MarRareqLogfrontier.activatedAugmentedSample_prior_singleton_eq

\item \emph{The cutoff probability and total variation.}
The retained baseline, treatment, and flag pattern has the same distribution
under both priors. In particular, for \(x<J\),
\[
V_x\sim\operatorname{Binomial}(n,b_0qH).
\]
Define
\[
k_{\mathrm{tail}}=K+1,\qquad
\nu_{\mathrm{act}}=nb_0qH=\frac{\eta K^2}{\ell}.
\]
For an integer \(v\ge0\), let its falling factorial be
\[
(v)_{k_{\mathrm{tail}}}
=
\begin{cases}
\displaystyle\prod_{t=0}^{k_{\mathrm{tail}}-1}(v-t),
&v\ge k_{\mathrm{tail}},\\
0,&v<k_{\mathrm{tail}}.
\end{cases}
\]
Expanding over ordered distinct record indices gives
\[
E[(V_x)_{k_{\mathrm{tail}}}]
=(n)_{k_{\mathrm{tail}}}(b_0qH)^{k_{\mathrm{tail}}}
\le\nu_{\mathrm{act}}^{k_{\mathrm{tail}}}.
\]
When \(k_{\mathrm{tail}}>n\), both falling-factorial moments are zero.
On \(V_x>K\), the falling factorial is at least \(k_{\mathrm{tail}}!\).
Therefore
\[
\widetilde Q_{\mathrm{act},\iota}(V_x>K)
\le
\frac{\nu_{\mathrm{act}}^{k_{\mathrm{tail}}}}
{k_{\mathrm{tail}}!}.
\]

The bounds \(K\le2\ell\) and \(K\le k_{\mathrm{tail}}\) imply
\[
K^2\le2\ell k_{\mathrm{tail}},
\qquad
\nu_{\mathrm{act}}\le2\eta k_{\mathrm{tail}}
\le e^{-31}k_{\mathrm{tail}},
\qquad
k_{\mathrm{tail}}\ge\ell,
\]
where \(2\le e\) gives the second inequality involving the exponential.
The exponential series gives
\(k_{\mathrm{tail}}^{k_{\mathrm{tail}}}/k_{\mathrm{tail}}!
\le e^{k_{\mathrm{tail}}}\), whence
\[
\frac{\nu_{\mathrm{act}}^{k_{\mathrm{tail}}}}
{k_{\mathrm{tail}}!}
\le
e^{-31k_{\mathrm{tail}}}
\frac{k_{\mathrm{tail}}^{k_{\mathrm{tail}}}}
{k_{\mathrm{tail}}!}
\le e^{-30k_{\mathrm{tail}}}
\le e^{-30\ell}.
\]
A union bound over the rare cells yields
\[
\widetilde Q_{\mathrm{act},\iota}(\mathcal G^c)
\le Je^{-30\ell}.
\]
The mass cap and elementary exponential bounds give
\[
J\le\frac{N\ell}{2\eta},\qquad
N\le e^\ell,\qquad \ell\le e^\ell,\qquad
N\ell\le e^{2\ell}.
\]
Consequently,
\[
\widetilde Q_{\mathrm{act},\iota}(\mathcal G^c)
\le
\frac{N\ell}{2\eta}e^{-30\ell}
\le
\frac{e^{2\ell}}{\eta}e^{-30\ell}
=e^{32-28\ell}.
\]
The threshold \(N\ge e^{(32+\log8)/28}\) implies
\[
\ell\ge\frac{32+\log8}{28},
\qquad
\widetilde Q_{\mathrm{act},\iota}(\mathcal G^c)\le\frac18.
\]

Equality of the atoms on \(\mathcal G\) implies equality of the probabilities
of \(\mathcal G\), and hence of its complement. For every augmented-sample
event \(\mathcal B\),
\[
\begin{aligned}
&\left|
\widetilde Q_{\mathrm{act},0}(\mathcal B)
-\widetilde Q_{\mathrm{act},1}(\mathcal B)
\right|\\
&\quad=
\left|
\widetilde Q_{\mathrm{act},0}(\mathcal B\cap\mathcal G^c)
-\widetilde Q_{\mathrm{act},1}(\mathcal B\cap\mathcal G^c)
\right|
\le\widetilde Q_{\mathrm{act},0}(\mathcal G^c).
\end{aligned}
\]
Applying this inequality to preimages under the forgetting map gives
\[
\operatorname{TV}(Q_{\mathrm{act},0},Q_{\mathrm{act},1})
\le
\operatorname{TV}(\widetilde Q_{\mathrm{act},0},
\widetilde Q_{\mathrm{act},1})
\le\frac18.
\]
% lean: CausalSmith.Stat.MarRareqLogfrontier.activatedAugmentedMixture_cutoff_budget

\item \emph{From center separation to squared-error risk.}
For \(\iota\in\{0,1\}\), define the prior center and its separation by
\[
\theta_{\mathrm{act},\iota}
=E_{\Pi_\iota^{\otimes d}}[\tau_{\mathrm{act}}(\mathbf Z)]
=Jb_0\int z^{-1}\,d\Pi_\iota(z),
\qquad
\Delta_{\mathrm{act}}
=|\theta_{\mathrm{act},1}-\theta_{\mathrm{act},0}|.
\]
The reciprocal-mean separation implies
\[
\Delta_{\mathrm{act}}\ge\frac{Jb_0}{12}.
\]
Since \(Z_x^{-1}\in[0,1]\), its variance satisfies
\[
\operatorname{Var}_{\Pi_\iota}(Z_x^{-1})
=
E_{\Pi_\iota}\!\left[(Z_x^{-1}-\tfrac12)^2\right]
-\left(E_{\Pi_\iota}[Z_x^{-1}]-\tfrac12\right)^2
\le\frac14.
\]
Independence of the latent coordinates therefore yields
\[
E_{\Pi_\iota^{\otimes d}}\!\left[
(\tau_{\mathrm{act}}(\mathbf Z)-\theta_{\mathrm{act},\iota})^2
\right]
=
b_0^2\sum_{x=0}^{J-1}
\operatorname{Var}_{\Pi_\iota}(Z_x^{-1})
\le\frac{Jb_0^2}{4}.
\]

Fix \(\mathsf T\in\mathcal T_n\), and define its maximal model risk and its
centered mixture risks by
\[
\mathcal W(\mathsf T)
=
\sup_{P\in\mathcal M_{n,d,q}}
E_P[(\mathsf T-\tau(P))^2],
\qquad
C_{\mathrm{act},\iota}(\mathsf T)
=
\int
(b_{\mathsf T}(o)-\theta_{\mathrm{act},\iota})^2\,
dQ_{\mathrm{act},\iota}(o).
\]
The bounded action and target make these risks finite. Class membership of
every prior support law gives
\[
\int
E_{P_{\mathrm{act},\mathbf z}}\!\left[
(\mathsf T-\tau_{\mathrm{act}}(\mathbf z))^2
\right]
d\Pi_\iota^{\otimes d}(\mathbf z)
\le\mathcal W(\mathsf T).
\]
Using the barycenter comparison from
\cref{obj:lem:randomized-kernel-barycenter} and
\[
(b_{\mathsf T}(o)-\theta_{\mathrm{act},\iota})^2
\le
2(b_{\mathsf T}(o)-\tau_{\mathrm{act}}(\mathbf z))^2
+
2(\tau_{\mathrm{act}}(\mathbf z)-\theta_{\mathrm{act},\iota})^2
\]
inside the mixture integral gives
\[
C_{\mathrm{act},\iota}(\mathsf T)
\le2\mathcal W(\mathsf T)+\frac{Jb_0^2}{2}.
\]

Define the nearer-center test
\[
\psi_{\mathrm{act}}(o)
=
\mathbf1\left\{
|b_{\mathsf T}(o)-\theta_{\mathrm{act},1}|
\le
|b_{\mathsf T}(o)-\theta_{\mathrm{act},0}|
\right\}.
\]
The triangle inequality shows that a testing error under center \(\iota\)
entails
\[
|b_{\mathsf T}(o)-\theta_{\mathrm{act},\iota}|
\ge\frac{\Delta_{\mathrm{act}}}{2}.
\]
By \cref{obj:lem:randomized-kernel-tv-comparison},
\[
Q_{\mathrm{act},0}(\psi_{\mathrm{act}}=1)
+
Q_{\mathrm{act},1}(\psi_{\mathrm{act}}=0)
\ge
1-\operatorname{TV}(Q_{\mathrm{act},0},Q_{\mathrm{act},1}).
\]
Bounding each error probability by its centered squared risk consequently
gives
\[
\frac{\Delta_{\mathrm{act}}^2}{8}
\left[
1-\operatorname{TV}(Q_{\mathrm{act},0},Q_{\mathrm{act},1})
\right]
\le
\max_{\iota\in\{0,1\}}C_{\mathrm{act},\iota}(\mathsf T).
\]
Together with the total variation and centered-risk bounds, this implies
\[
\frac{7\Delta_{\mathrm{act}}^2}{128}
-\frac{Jb_0^2}{4}
\le\mathcal W(\mathsf T).
\]
Now \(J\ge2048\) and \(\Delta_{\mathrm{act}}\ge Jb_0/12\), so
\[
\mathcal W(\mathsf T)
\ge
(Jb_0)^2
\left(\frac7{18432}-\frac1{4J}\right)
\ge\frac{(Jb_0)^2}{4096}.
\]
Taking the infimum over \(\mathsf T\) gives the calibrated fuzzy bound
\[
\frac{(Jb_0)^2}{4096}\le\mathfrak R_{n,d,q}.
\]
% lean: CausalSmith.Stat.MarRareqLogfrontier.activatedAugmentedMixture_minimaxRisk_calibrated

\item \emph{Converting rare-cell mass to the dimensional contribution.}
Still suppose \(N\ge N_{\mathrm{act}}\) and \(J\ge2048\). We have
\[
b_0=\frac{\eta}{N\ell}\le\frac14,
\qquad
0<\eta=e^{-32}\le\frac12,
\qquad d\ge2.
\]
The first inequality uses \(N\ge4\eta\) and \(\ell\ge1\); the second follows
from \(e^{32}\ge1+32\).

If \(d-1\le\lfloor(2b_0)^{-1}\rfloor\), then
\[
J=d-1\ge\frac d2,\qquad Jb_0\ge\frac{db_0}{2}.
\]
Otherwise the floor determines \(J\), and
\[
Jb_0
=\left\lfloor\frac1{2b_0}\right\rfloor b_0
\ge\frac12-b_0\ge\frac14.
\]
Thus both branches give
\[
\min\left\{\frac{db_0}{2},\frac14\right\}\le Jb_0.
\]
Moreover,
\[
\frac{\eta}{24}\min\{1,u_{\mathrm{act}}\}
\le\frac{\eta u_{\mathrm{act}}}{24}
=\frac{db_0}{24},
\qquad
\frac{\eta}{24}\min\{1,u_{\mathrm{act}}\}
\le\frac{\eta}{24}\le\frac1{48}.
\]
It follows that
\[
\frac{\eta}{24}\min\{1,u_{\mathrm{act}}\}
\le
\frac1{12}\min\left\{\frac{db_0}{2},\frac14\right\}
\le\frac{Jb_0}{12},
\]
and hence
\[
\frac{\eta}{2}\min\{1,u_{\mathrm{act}}\}\le Jb_0.
\]
For \(u_{\mathrm{act}}\ge0\), considering \(u_{\mathrm{act}}\le1\) and
\(u_{\mathrm{act}}\ge1\) gives
\[
\bigl(\min\{1,u_{\mathrm{act}}\}\bigr)^2
=\min\{1,u_{\mathrm{act}}^2\}.
\]
Squaring the mass bound therefore yields
\[
\frac{\eta^2}{4}\min\{1,u_{\mathrm{act}}^2\}
\le(Jb_0)^2.
\]
The choice of \(\underline c_{\mathrm{act}}\) and the calibrated fuzzy bound
now imply
\[
\underline c_{\mathrm{act}}\min\{1,b_{\mathrm{dim}}\}
\le
\frac{\eta^2}{16384}\min\{1,u_{\mathrm{act}}^2\}
\le\frac{(Jb_0)^2}{4096}
\le\mathfrak R_{n,d,q}.
\]
% lean: CausalSmith.Stat.MarRareqLogfrontier.interval_prior_segment_width_lower

\item \emph{Large effective size with fewer rare cells.}
Suppose instead that \(N\ge N_{\mathrm{act}}\) and \(J<2048\).
The capacity bound established above gives
\[
\left\lfloor\frac1{2b_0}\right\rfloor\ge2048.
\]
The definition of \(J\) therefore forces
\[
d-1<2048,\qquad d\le2048.
\]
This includes \(J=0\), for which the same argument gives \(d=1\).

Since \(N\ge1\) and \(\ell\ge1\),
\[
\sqrt N\le N\le N\ell.
\]
Consequently,
\[
\min\{1,u_{\mathrm{act}}\}
\le\frac{d}{N\ell}
\le\frac{2048}{\sqrt N},
\qquad
\min\{1,u_{\mathrm{act}}^2\}
\le\frac{2048^2}{N}.
\]
Also,
\[
\min\{1,N^{-1}\}=N^{-1},
\qquad
\underline c_{\mathrm{act}}\,2048^2\le\frac1{256}.
\]
The numerical one-cell estimate thus gives
\[
\underline c_{\mathrm{act}}\min\{1,b_{\mathrm{dim}}\}
\le
\frac{\underline c_{\mathrm{act}}\,2048^2}{N}
\le\frac1{256N}
\le\mathfrak R_{n,d,q}.
\]
% lean: CausalSmith.Stat.MarRareqLogfrontier.activated_fuzzy_lower

\item \emph{Small effective size.}
In the remaining effective-size branch, \(N<N_{\mathrm{act}}\).
Because \(N>0\) and \(N_{\mathrm{act}}\ge1\),
\[
N_{\mathrm{act}}^{-1}\le1,\qquad
N_{\mathrm{act}}^{-1}\le N^{-1},
\qquad
N_{\mathrm{act}}^{-1}\le\min\{1,N^{-1}\}.
\]
The definition of \(\underline c_{\mathrm{act}}\) and the numerical
one-cell estimate give
\[
\underline c_{\mathrm{act}}\min\{1,b_{\mathrm{dim}}\}
\le\underline c_{\mathrm{act}}
\le\frac1{256N_{\mathrm{act}}}
\le\frac1{256}\min\{1,N^{-1}\}
\le\mathfrak R_{n,d,q}.
\]
The three effective-size and rare-count branches establish the complementary
bound whenever \(\sqrt N\le d\).
% lean: CausalSmith.Stat.MarRareqLogfrontier.activated_fuzzy_lower

\item \emph{Combining the bounds when \(d\ge\sqrt N\).}
First suppose \(\sqrt N\le d\). Both component bounds apply, and their
nonnegative capped factors give
\[
\underline c_*\min\{1,a_{\mathrm{eff}}\}
\le
\underline c_{\mathrm{one}}\min\{1,a_{\mathrm{eff}}\}
\le\mathfrak R_{n,d,q},
\]
\[
\underline c_*\min\{1,b_{\mathrm{dim}}\}
\le
\underline c_{\mathrm{act}}\min\{1,b_{\mathrm{dim}}\}
\le\mathfrak R_{n,d,q}.
\]

For the nonnegative quantities \(a_{\mathrm{eff}},b_{\mathrm{dim}}\),
\[
\min\{1,a_{\mathrm{eff}}+b_{\mathrm{dim}}\}
\le
\min\{1,a_{\mathrm{eff}}\}+\min\{1,b_{\mathrm{dim}}\}.
\]
Indeed, if both quantities are at most one, the right-hand side is their sum;
if either is at least one, the right-hand side is at least one.
Each capped quantity is at most their maximum, so
\[
\min\{1,a_{\mathrm{eff}}+b_{\mathrm{dim}}\}
\le
2\max\{\min\{1,a_{\mathrm{eff}}\},\min\{1,b_{\mathrm{dim}}\}\},
\]
while the component risk bounds imply
\[
\underline c_*
\max\{\min\{1,a_{\mathrm{eff}}\},\min\{1,b_{\mathrm{dim}}\}\}
\le\mathfrak R_{n,d,q}.
\]
Since \(\mathfrak R_{n,d,q}\ge0\), we conclude that
\[
\underline c\,r_{n,d,q}
\le
\frac{\underline c_*}{2}
\max\{\min\{1,a_{\mathrm{eff}}\},\min\{1,b_{\mathrm{dim}}\}\}
\le\frac12\mathfrak R_{n,d,q}
\le\mathfrak R_{n,d,q}.
\]
% lean: CausalSmith.Stat.MarRareqLogfrontier.full_experiment_lower

\item \emph{Combining the bounds when \(d<\sqrt N\).}
Suppose \(\sqrt N>d\). Since \(d\ge0\),
\[
d^2\le(\sqrt N)^2=N.
\]
Also \(\ell\ge1\), so
\[
\ell^2\ge1,\qquad
(N\ell)^2>0,\qquad
N^{-1}(N\ell)^2=N\ell^2.
\]
Therefore
\[
b_{\mathrm{dim}}
=\frac{d^2}{(N\ell)^2}
\le\frac{N\ell^2}{(N\ell)^2}
=N^{-1}
=a_{\mathrm{eff}},
\]
and monotonicity of the minimum gives
\[
\min\{1,b_{\mathrm{dim}}\}
\le\min\{1,a_{\mathrm{eff}}\}.
\]
Applying the capped-sum inequality established above,
\[
r_{n,d,q}
\le
\min\{1,a_{\mathrm{eff}}\}+\min\{1,b_{\mathrm{dim}}\}
\le2\min\{1,a_{\mathrm{eff}}\}.
\]
Finally, the one-cell contribution and
\(\underline c_*\le\underline c_{\mathrm{one}}\) yield
\[
\underline c\,r_{n,d,q}
\le
\frac{\underline c_*}{2}\min\{1,a_{\mathrm{eff}}\}
\le
\underline c_{\mathrm{one}}\min\{1,a_{\mathrm{eff}}\}
\le\mathfrak R_{n,d,q}.
\]
The two dimension branches cover all admissible parameters.
% lean: CausalSmith.Stat.MarRareqLogfrontier.full_experiment_lower
\end{enumerate}
\end{proof}
